\subsection{Certified absorption of disk-base annuli}
\label{marked-cut:disk-absorption}

The next implementation checkpoint removes a specific obstruction left by
whole-circle interfaces. A general annulus is still outside the disk-gluing
criterion, but a complete vertical annulus of a disk-base bundle determines
an actual two-handle. The implementation now constructs its source attaching
curve, authenticates it independently, and includes that handle in the
geometric cancellation proof. This implements the absorption proposition of
report~85; it does not interpret an arbitrary presentation row
as an embedded attaching curve.

\begin{proposition}[Augmented source core]
\label{marked-cut:absorbed-core}
Suppose a source-certified bundle component has a disk base and its entire
base boundary is attached to the core. It can be absorbed as one two-handle,
without adding zero-, one-, or three-handles. Its raw intersection count with
one-handle $h$ is the number of its core-one ports owned by $h$, counting
all occurrences without signs or cancellation. Any collection of these
components can be absorbed simultaneously. The added handles and their total
raw occurrences number at most $P$, the number of source ports.
\end{proposition}
\begin{proof}
The bundle over a disk is $D^2\times I$, and its complete attaching region
is $\partial D^2\times I$. This is the defining attaching region of a
three-dimensional two-handle. The two horizontal disks remain free boundary.
Take the midpoint of every interval fibre along the attaching annulus.
All source transition maps, including fibre reversals, preserve the midpoint,
so these local sections define an embedded attaching circle.

A core-zero port contributes an arc within its named zero-handle boundary.
A core-one port is a face-sector side of an edge-cap polygon. Its midpoint
arc traverses the lateral boundary of the corresponding source one-handle
from one end to the other, meeting its cocore once. Thus each such port
contributes one unsigned occurrence. Two occurrences with opposite signs
still contribute two intersections.

The complete lower attachments of the complementary source handles remain
in the core. The product contacts occupy their recorded boundary regions,
disjoint from the original two-handle attaching regions after compatible
corner rounding. Consequently the new attaching circles can be placed at
the same two-handle stage as the old ones. Distinct bundle components have
disjoint attaching annuli. Absorption therefore changes neither the embedded
cut manifold nor the interfaces of the remaining bundles. Each absorbed
component uses at least one port, and distinct components use disjoint
ports, proving both bounds.
\end{proof}

\paragraph{Face flags authenticate the actual path.}
Write the two source ends of one-handle $h$ as
$(t_h,f_h)$ and $(u_h,g_h)$. For a core-one port inspect its two cap-edge
endpoint flags. Their order must match these two flags or their reverse.
Call the resulting local sign $\delta_p\in\{+1,-1\}$. If the canonical
boundary word traverses that port with sign $s_p$, the corresponding raw
one-handle crossing has sign $s_p\delta_p$. Its initial and terminal
zero-handles are read from the appropriate ends of $h$.

Both tetrahedron and face are essential. In a self-paired tetrahedron the
one-handle can have equal tetrahedron labels and even equal zero-handle
owners at its ends, while its two face flags still differ. Comparing only
tetrahedron labels would lose this distinction. Every core-zero port stays
at its owner vertex. The producer and independent checker require the end
of every oriented port path to equal the beginning of the next, including
the final-to-initial join. Deleting the stationary portions then gives the
alternating zero/one cycle used by the existing core checker. The signs
are retained as chart evidence; the cancellation matrix remains unsigned.
A circle with no one-handle crossings contributes a zero row and cannot
silently disappear from the required cancellation schedule.

\paragraph{Eligibility and independent replay.}
The adapter requires a single base owner, signature $(\chi,b)=(1,1)$,
orientability, a product bundle, and multiplicity one. The circle must
exhaust that component's port set, and no component may be absorbed twice.
These facts are obtained from the source-verified base profiles and ordered
interface, rather than from the proposed absorption record. Because the
base has exactly one boundary component, a complete circle uses its whole
vertical boundary. A component with a nondisk base remains unresolved.

The new whole-cut certificate schema is
\path{normal-cut-handlebodies-v2}. Its absorption record names the source
attachment, base profile, appended two-handle, and every directed crossing.
The independent module \path{disk_bundle_absorption_verify} reconstructs
these data from authenticated polygons and endpoint flags. It imports no
producer. Only after exact comparison does it return an augmented geometric
core to the existing independent cancellation checker. Every added row must
be eliminated; the genus calculation uses the surviving one-handles of
this augmented core. Version-one certificates remain replayable, and cuts
requiring no absorption continue to use version one.

\paragraph{The whole-cut census after absorption.}
Remove absorbed bases and their annular attachments from the bundle/core
graph, replacing the old core by the augmented core. All remaining bases
must be bordered and all their attachments must be proper intervals.
The disk-gluing theorem then gives the same genus formula
\[
 g(X)=\sum_{v\in V(X)}g(v)+|E(X)|-|V(X)|+1.
\]
An absorbed disk contributes one to the core Euler characteristic and
removes a bundle of Euler characteristic one; its attaching annulus has
Euler characteristic zero. Hence the original source-cut Euler identity
remains the correct final check. Unported bundle multiplicities are counted
as before. This proves the complete histogram only when all augmented core
cancellations and all remaining disk gluings succeed.

\paragraph{Cost and scope.}
The adapter adds $O(P)$ records and $O(P\log(T+P))$ bits of finite chart data;
it never expands the binary normal multiplicities. Constructing the records
is linear in the finite source data, while the checker also sorts port sets
in $O(P\log P)$ comparisons. The augmented core still has $O(T)$ handles and
raw occurrences, so the existing polynomial saturated-cancellation bounds
apply unchanged in order. The six weighted orbit queries and their shared
budgets are unchanged. Failure of a cancellation schedule or rejection of a
nondisk annulus carries no negative topology conclusion. This advances the
supplied-cut terminal operator, but does not provide useful-surface discovery,
arbitrary boundary-pattern transport, or a global quasi-polynomial hierarchy
bound.

\paragraph{Source-cut audit.}
The refreshed 106-case audit certifies 63 whole cuts, compared with 30 before
absorption. Every new success is among the 33 algebraically eliminable
candidates identified by the earlier probe, and all previous successes are
preserved. All 63 source certificates replay independently. Regina's actual
cut construction and component recognition finish for 99 cases and agree
exactly with every one of the 56 native successes within the expansion
limit. Seven large-coordinate successes exceed the 4,096-disk literal cap
and are independently replayed without expansion. No oracle times out.
The source pins and per-case results are retained in
\path{data/disk-absorption-audit.json}.

The 43 remaining refusals comprise 26 nondisk annular cases, seven stalled
augmented cancellations (three detected by the square leaf obstruction),
and ten cores with three-handles. These categories are not perfectly
comparable with the earlier refusal ordering: the boundary and absorption
checks now precede cancellation, so four previously stalled core cases
are classified first by their nondisk annuli. They are not lost successes.

The focused checks cover actual same-tetrahedron and same-zero-handle
self-pairings, source relabellings, signed endpoint charts, malformed and
omitted absorption records, an omitted cancellation, disabled producers,
exact hexadecimal transport, resource caps, nondisk refusals and stalled
augmented cores. Multiplying the two-tetrahedron annular cut by $2^{1024}$
produces and independently verifies exactly $2^{1024}+1$ solid-torus
components without enumerating them. The tests also retain the pre-existing
raw cancellation and disk-gluing checks. These are source-bound whole-cut
certificates, not an enlargement of the complete useful-surface search.

